\documentclass{article}
\usepackage[a4paper, total={6.5in, 10in}]{geometry}
\setlength{\parindent}{0pt}
\usepackage{tcolorbox}
\tcbset{colback=white!94!black, colframe=black!60!white, coltitle = white, boxrule=0.8pt, arc=2mm}
\newtcolorbox{boxs}[2][]{title= #2,#1}
\usepackage{datetime2}
\usepackage[british]{babel}
\usepackage{amsfonts}
\usepackage{amsmath}

\title{\textbf{Linear Algebra HW 9}}
\author{Robert O'Connell}
\date{\today}
\begin{document}
\maketitle
\vspace{5mm}
\subsection*{Question 1}
Let \(\vec{w_1} = \vec{v_1}\)

\begin{align*}
    \vec{w_2} &= \vec{v_2} - \frac{\langle \vec{w_1}, \vec{v_2} \rangle}{\langle \vec{w_1}, \vec{w_1} \rangle} \vec{w_1} \\
    &= \begin{bmatrix}
        1 \\ 0 \\ 1 
    \end{bmatrix} - \frac{2}{6} \begin{bmatrix}
        1 \\ 2 \\ 1 
    \end{bmatrix} \\
    &= \begin{bmatrix}
        \frac{2}{3} \\ -\frac{2}{3} \\ \frac{2}{3}
    \end{bmatrix}
\end{align*}

\begin{align*}
\vec{w_3} &= \vec{v_3} -  \frac{\langle \vec{w_1}, \vec{v_3} \rangle}{\langle \vec{w_1}, \vec{w_1} \rangle} \vec{w_1} -  \frac{\langle \vec{w_2}, \vec{v_3} \rangle}{\langle \vec{w_2}, \vec{w_2} \rangle} \vec{w_2} \\
&= \begin{bmatrix}
    0 \\ 1 \\ 1
\end{bmatrix} - \frac{3}{6}\begin{bmatrix}
    1 \\ 2 \\ 1
\end{bmatrix} - (0)(\dots) \\
&= \begin{bmatrix}
    -\frac{1}{2} \\ 0 \\ \frac{1}{2}
\end{bmatrix}
\end{align*}

\subsection*{Question 2}
Let the basis \(\{1,x,x^2,x^3\}\) have entries \(\vec{v_i}, \quad 1 \leq i \leq 4\)

Let \(\vec{w_1} = \vec{v_1}\)

\textbf{Note:}
If \(f(x) = x^i\), \(i\) odd, \(\int_{-a}^{a} f(x) dx = 0\)

\begin{align*}
    \vec{w_2} &= \vec{v_2}  - \frac{\langle \vec{w_1}, \vec{v_2} \rangle}{\langle \vec{w_1}, \vec{w_1} \rangle} \vec{w_1} \\
    &= x - \frac{\int_{-2}^{2} x dx}{\int_{-2}^{2} dx} \\
    &= x
\end{align*}

\begin{align*}
    \vec{w_3} &= x^2 - \frac{\int_{-2}^{2}x^2 dx}{\int_{-2}^{2}dx} - \frac{\int_{-2}^{2}x^3 dx}{\int_{-2}^{2} x^2 dx}x \\
    &= x^2 - \frac{\frac{16}{3}}{4} \\
    &= x^2 - \frac{4}{3}
\end{align*}

\begin{align*}
    \vec{w_4} &= x^3 - \frac{\int_{-2}^{2}x^3dx}{\int_{-2}^{2} dx} - \frac{\int_{-2}^{2}x^4 dx}{\int_{-2}^{2}x^2 dx}x - \frac{\int_{-2}^{2}(x^2 - \frac{4}{3})(x^3) dx}{\int_{-2}^{2}x^4 dx}x^2 \\
    &= x^3 - \frac{\frac{64}{5}}{\frac{16}{3}}x \\
    &= x^3 - \frac{12}{5}x
\end{align*}

Hence \(\left\{ \vec{w_1}, \vec{w_2}, \vec{w_3}, \vec{w_4} \right\}\) form an orthogonal basis for \(\mathbb{P}_2[x]\) with the given inner product

To get an orthonormal basis, we simply divide by the norm of the vectors w.r.t. to the inner product

\[
\langle \vec{w_1}, \vec{w_1} \rangle = \int_{-2}^{2} dx = 4
\]
\[
\langle \vec{w_2}, \vec{w_2} \rangle = \int_{-2}^{2} x^2 dx = \frac{16}{3}
\]
\[
\langle \vec{w_3}, \vec{w_2} \rangle = \int_{-2}^{2} x^4 dx = \frac{64}{5}
\]
\[
\langle \vec{w_4}, \vec{w_4} \rangle = \int_{-2}^{2} x^6 dx = \frac{256}{7}
\]

Hence \(\left\{ \frac{1}{4}, \frac{3x}{16}, \frac{5(x^2 - \frac{4}{3})}{64}, \frac{7(x^3 - \frac{12}{5}x)}{256}\right\}\), forms an orthonormal basis











\subsection*{Question 3}
\subsubsection*{(a)}
(i)

Let \(A,B,C \in M_{2\times 2}(\mathbb{R})\), \(\lambda \in \mathbb{R}\)

Checking for linearity in the first variable, 
\begin{align*}
    \langle \lambda A + C , B\rangle &= \mathrm{tr}((\lambda A + C)B) \\
    &= \mathrm{tr}(\lambda AB + CB) \\
    &= \lambda \mathrm{tr}(AB) + \mathrm{tr}(CB)
\end{align*}
Hence it is linear in the first variable

Checking for linearity in the second variable, 
\begin{align*}
    \langle  A , \lambda B + C\rangle &= \mathrm{tr}(A(\lambda B + C)) \\
    &= \mathrm{tr}(\lambda AB + AC) \\
    &= \lambda \mathrm{tr}(AB) + \mathrm{tr}(AC)
\end{align*}
Hence it is linear in the second variable

Hence it is bilnear
\\

(ii)

\begin{align*}
    \langle A,B \rangle &= \mathrm{tr}(AB)
\end{align*}

With \(A\) having \(ij\) entries \(a_{ij}\), and similar for \(B\), we get 
\[
AB = \begin{bmatrix}
    a_{11}b_{11} + a_{12}b_{21} & \dots \\
    \dots & a_{21}b_{12} + a_{22}b_{22}
\end{bmatrix}
\]
\[
\Rightarrow \mathrm{tr}(AB) = a_{11}b_{11} + a_{12}b_{21} + a_{21}b_{12} + a_{22}b_{22}
\]
From this we see that \(\mathrm{tr}(AB) = \mathrm{tr}(BA)\)

Then
\[
\langle A,B \rangle = \langle B,A \rangle
\]
Hence it is symmetric
\\

(iii)
\[
\langle A,A \rangle = \mathrm{tr}(A^2)
\]
This is not always greater or equal to zero

Consider \(A = \begin{bmatrix}
    0 & 1 \\
    -1 & 0  
\end{bmatrix}\), \(A^2 = \begin{bmatrix}
    -1 & 0 \\
    0 & -1 
\end{bmatrix}\)

Hence quite clearly \(\mathrm{tr}(A^2) < 0\)





\subsubsection*{(b)}
(i)

Let \(A,B,C \in M_{2\times 2}(\mathbb{R})\), \(\lambda \in \mathbb{R}\)

Checking for linearity in the first variable, 
\begin{align*}
    \langle \lambda A + C , B\rangle &= \mathrm{tr}((\lambda A + C)^T B) \\
    &= \mathrm{tr}(\lambda A^TB + C^T B) \\
    &= \lambda \mathrm{tr}(A^T B) + \mathrm{tr}(C^TB)
\end{align*}
Hence it is linear in the first variable

Checking for linearity in the second variable, 
\begin{align*}
    \langle  A , \lambda B + C\rangle &= \mathrm{tr}(A^T(\lambda B + C)) \\
    &= \mathrm{tr}(\lambda A^TB + A^TC) \\
    &= \lambda \mathrm{tr}(A^TB) + \mathrm{tr}(A^TC)
\end{align*}
Hence it is linear in the second variable

Hence it is bilnear
\\

(ii)

\begin{align*}
    \langle A,B \rangle &= \mathrm{tr}(A^TB)
\end{align*}

With \(A\) having \(ij\) entries \(a_{ij}\), and similar for \(B\), we get 
\[
A^T B = \begin{bmatrix}
    a_{11}b_{11} + a_{21}b_{21} & \dots \\
    \dots & a_{12}b_{12} + a_{22}b_{22}
\end{bmatrix}
\]
\[
\Rightarrow \mathrm{tr}(A^TB) = a_{11}b_{11} + a_{21}b_{21} + a_{12}b_{12} + a_{22}b_{22}
\]
From this we see that \(\mathrm{tr}(A^TB) = \mathrm{tr}(B^TA)\)

Then
\[
\langle A,B \rangle = \langle B,A \rangle
\]
Hence it is symmetric
\\

(iii)
\[
\langle A,A \rangle = \mathrm{tr}(A^TA)
\]
With \(ij\) entries as before we get
\[
A^TA = \begin{bmatrix}
    a_{11}^2 + a_{12}^2 & \dots \\
    \dots & a_{21}^2 + a_{22}^2
\end{bmatrix}
\]
Hence \(\mathrm{tr}(A^TA) \geq 0\), with equality \(\iff\) \(a_{ij} = 0 , \quad \forall \ 1 \leq i,j \leq 2\)

\subsubsection*{(c)}
(i)

Let \(A,B,C \in M_{2\times 2}(\mathbb{R})\), \(\lambda \in \mathbb{R}\)

Checking for linearity in the first variable
\begin{align*}
    \langle \lambda A + C, B \rangle &= \det(\lambda A + C + B) \\
\end{align*}

Since the determinant does not preserve addition, this is not bilnear
\\

(ii)
\[
\langle A,B \rangle = \det(A+B) = \det(B+A) = \langle B,A \rangle
\]
Hence it is symmetric
\\

(iii)
\[
\langle A,A \rangle = \det(A+A) = 2\det(A)
\]
But we know the determinant of a matrix can be negative

Thus it is not postive definite










\subsection*{Question 4}
Since \(A\) is real symmetric matrix there exists an orthogonal matrix \(B\) s.t.
\(B^{-1} AB \) is diagonal

\begin{align*}
    \chi_A(\lambda) &= (2-\lambda)((3-\lambda)^2 - 1) \\
    &= (2-\lambda)(\lambda^2 -6\lambda + 8) \\
    &= (2-\lambda)^2(\lambda -4)
\end{align*}
Thus there are two eigenvalues \(\lambda = 2\) and \(\lambda = 4\)

\(\mathcal{N}(A-4I)\):
\[
\begin{bmatrix}
    -1 & 1 & 0 & | & 0 \\
    1 & -1 & 0 & | & 0 \\
    0 & 0 & 2 & | & 0 
\end{bmatrix}
\]
\[
\Rightarrow \begin{bmatrix}
    1 & -1 & 0 & | & 0 \\
    0 & 0 & 0 & | & 0 \\
    0 & 0 & 0 & | & 0 \\
\end{bmatrix}
\]
\[
\Rightarrow V(4) = \left\{ \begin{bmatrix}
t \\ t \\ 0
\end{bmatrix} | t \in \mathbb{R}  \right\}
\]

\(\mathcal{N}(A-2I)\):
\[
\begin{bmatrix}
    1 & 1 & 0 & | & 0 \\
    1 & 1 & 0 & | & 0 \\
    0 & 0 & 0 & | & 0 
\end{bmatrix}
\]
\[
\Rightarrow \begin{bmatrix}
    1 & 1 & 0 & | & 0 \\
    0 & 0 & 0 & | & 0 \\
    0 & 0 & 0 & | & 0 \\
\end{bmatrix}
\]
\[
\Rightarrow V(2) = \left\{ \begin{bmatrix}
t \\ -t \\ s
\end{bmatrix} | t,s \in \mathbb{R}  \right\}
\]

Since eigenvectors associated with different eigenvalues of a symmetric matrix are always
orthogonal, we only need to be careful with the two eigenvectors associated with \(\lambda = 2\)

Choosing them to be \(\begin{bmatrix}
    1 \\ -1 \\ 0
\end{bmatrix}\) and \(\begin{bmatrix}
    0 \\ 0 \\ 1
\end{bmatrix}\), with the eigenvector associated with \(\lambda = 4 \) being \(\begin{bmatrix}
    1 \\ 1 \\ 0
\end{bmatrix}\), we get 

\[
B = \begin{bmatrix}
    1 & 1 & 0 \\
    1 & -1 & 0 \\
    0 & 0 & 1
\end{bmatrix}
\]


\subsection*{Question 5}
The matrix \(A\) defines a bilinear form by \(\langle \vec{x} , \vec{y} \rangle = \vec{x}^T A \vec{y}, \quad \vec{x}, \vec{y} \in \mathbb{R}^3\)

Let \(\vec{x}\) have components \(x_1, x_2, x_3\), similarly for \(\vec{y}\)
\begin{align*}
    \langle \vec{x}, \vec{y}\rangle &= \begin{bmatrix}
        x_1 a + x_2 + x_3 & x_1  + x_2 a + x_3 & x_1 + x_2 + x_3 a
    \end{bmatrix} \vec{y} \\
    &= y_1(x_1 a + x_2 + x_3) + y_2(x_1  + x_2 a + x_3) + y_3 (x_1 + x_2 + x_3 a) 
\end{align*}

To check for postive defitiness, we compute \(\langle \vec{x}, \vec{x} \rangle\)

\begin{align*}
    \langle \vec{x}, \vec{x}\rangle &= x_1(x_1 a + x_2 + x_3) + x_2(x_1  + x_2 a + x_3) + x_3 (x_1 + x_2 + x_3 a)  \\
    &= a (x_1^2 + x_2^2 + x_3^2) + 2x_1x_3 + 2x_1x_2 + 2x_2x_3
\end{align*}

Since \((x_1 + x_2 + x_3)^2 = x_1^2+x_2^2+x_3^2 + 2(x_1x_2 + x_1x_3 + x_2x_3)\) 

i.e. \(2(x_1x_2 + x_1x_3 + x_2x_3) = (x_1 + x_2 + x_3)^2 - (x_1^2+x_2^2+x_3^2 ) \)

we have
\begin{align*}
    \langle \vec{x}, \vec{x}\rangle &= (a-1) (x_1^2 + x_2^2 + x_3^2) + (x_1+x_2+x_3)^2
\end{align*}

Thus we see that this is greater than equal to zero for \(a \geq 1\)

Hence \(A\) is postive definite for \(a \geq 1\)

\end{document}